\section{Introduction}
\label{sec:intro}

Selective state-space models (SSMs) of the Mamba family have rapidly become default sequence backbones, and a common narrative has formed around \emph{why} they work: input-dependent selectivity---the ability to modulate state dynamics per token---is credited with filtering noise, adapting to distribution shift, and outperforming time-invariant predecessors such as S4D~\cite{gu2023mamba}. \cite{qu2024mambasurvey,patro2024mamba360} This belief motivates a growing family of ``selective'' architectures, yet recent evidence questions it: a systematic comparison on 59 clean multivariate time-series classification benchmarks (MONSTER and UEA) finds the \emph{simpler, time-invariant} S4D consistently ahead of Mamba-family variants in accuracy and efficiency~\cite{saadatmand2026s4d}, and theoretical analyses identify expressivity regimes where Mamba is provably weaker~\cite{yu2025blockbiased}. Which account is right---and, more importantly, \emph{what exactly should be credited} when a Mamba-style block does prove robust? This paper addresses the second question for one block and one testbed, with a controlled ablation chain whose decision thresholds were fixed before each grid was run. The testbed is bearing vibration diagnosis under operating-condition drift and matched noise, in which the robustness gap is large (12--28$\pp$ between S4D and the frozen Mamba-3 blocks in the independent-noise results), the evaluation labels are identical across all cells of each transition (hash-verified), and the block can be dissected.

Failure to identify the source of an empirical gain is a recurring weakness of machine-learning scholarship~\cite{lipton2018troubling}. Our approach is deliberately \emph{attributional} rather than \emph{methodological}: we propose no new architecture and claim no state-of-the-art. We take a regime in which a Mamba-style block shows a large, stable robustness advantage over a time-invariant SSM, and ask which ingredient earns it, in three steps with adjudication thresholds fixed before each grid was run:

\begin{enumerate}
\item \textbf{Freeze the scan.} The input-dependent parts of $\Delta$, $A$, $B$, $C$ are replaced by learned constants; in a second arm the trapezoidal weight and rotation angles are as well, so that the scan's dynamics and read/write coefficients no longer depend on the input. The value stream and the gate input remain input-dependent (Section~\ref{sec:protocol:freeze}).
\item \textbf{Remove the gate.} On the frozen base the multiplicative SiLU gate is removed, with the other components and the budget unchanged.
\item \textbf{Replace or transplant the gate.} In place, the gate is replaced by an additive branch of identical size and initial weights; across hosts, it is grafted onto plain S4D with an additive control.
\end{enumerate}

Steps 1 and 2 and the in-place replacement were tested with training and evaluation noise drawn independently (Section~\ref{sec:results}); the graft onto S4D, the per-class analysis, the stem and $A$ ablations and the disjoint-bearing Paderborn setting were run only under an earlier protocol in which the two noise streams were shared by window index, and are reported in \ref{app:earlier} as descriptive evidence (Table~\ref{tab:coverage}). The outcome is a conditional pattern. Making the scan input-independent does not remove the block's advantage over S4D, on either transition ($+21.3$, $+23.3$, $+20.3$\,pp and $+27.5$, $+26.7$, $+21.1$\,pp for the fully frozen block). Removing the gate costs $+9.4$, $+15.7$, $+27.0$\,pp in the partially frozen block on the original transition, resolved only at $-6$\,dB, and 19.3, 22.7, 22.0\,pp on the reverse transition, resolved at every level. Only in the fully frozen block on the original transition does the gate-removal cost grow significantly with noise depth: without its gate that block still exceeds S4D at $0$ and $-2$\,dB but not at $-6$\,dB. An additive branch of identical size does not recover the gate's effect at $-6$\,dB in either host, although the pre-specified rescue test is not assessable. On Paderborn bearings the gate-removal cost replicates when training and evaluation share bearings.

\textbf{Contributions.} (i) A layered freezing experiment on a Mamba-3 block, down to a scan with input-independent dynamics and read/write coefficients (verified at the kernel interface), in which the advantage over S4D survives on both transitions of one dataset under independent training and evaluation noise. (ii) Same-host tests of the output gate (removal and in-place additive replacement) in a partially and a fully frozen block, with a descriptive analysis of how the gate's contribution changes with noise depth. As supporting material, the study reports a coverage table by protocol, the correction of an initial noise-protocol defect, and a released pre-specification and review trail, including the incidents in which automated adjudication departed from the pre-specified criteria (Appendix~C).
